\section{Analisis de resultados}

A partir de las mediciones presentadas en la sección de resultados, es posible efectuar las siguientes observaciones para cada tipo de carga ensayada.

\begin{itemize}
    \item \textbf{Carga lineal:} La relación de transformación se mantiene próxima a 2 ($V_2/V_1 \approx 220/110$) en todos los niveles de carga ensayados, lo cual es coherente con el ajuste nominal del transformador. Se aprecia una ligera caída de la tensión del secundario al aumentar la carga (de 220 V en 200 W a 215 V en 800 W y 1 kW), comportamiento esperado de la regulación de tensión de un transformador real bajo carga. No obstante, las lecturas de potencia presentan inconsistencias: la potencia activa del primario ($P_1$) resulta menor que la del secundario ($P_2$) en todos los niveles (por ejemplo, $P_1 = 304$ W frente a $P_2 = 864$ W a 1 kW), lo cual contradice el balance de potencias de una máquina pasiva con pérdidas internas. Adicionalmente, el factor de potencia calculado en el primario ($\approx 0,21$ a 1 kW) difiere notablemente del secundario ($\approx 0,98$), pese a tratarse de una carga resistiva de factor de potencia unitario. Cabe resaltar que el factor de potencia del secundario es muy próximo a la unidad (0,98), lo cual es coherente con una carga de carácter resistivo; la pequeña desviación respecto de 1, junto con la discrepancia del primario, se atribuye a errores de medición e incertidumbre de los instrumentos (rango de lectura del vatímetro), más que a un comportamiento físico real de la carga. Estas discrepancias sugieren errores de rango o de conexión del vatímetro en alguno de los lados del transformador. La curva de voltaje vs corriente compilada y las formas de onda senoidales capturadas confirman el comportamiento lineal esperado en este ensayo: al tratarse de una carga lineal, la forma de onda es lineal (senoidal) y el desfase entre voltaje y corriente es nulo. El pequeño desfase que se aprecia en las figuras puede atribuirse a los elementos inductivos del transformador (reactancia de dispersión e inductancia de magnetización), que introducen un ligero retraso de la corriente respecto a la tensión.

    \item \textbf{Carga capacitiva:} Con $V_1 = 68,5$ V se obtuvo $V_2 = 126,4$ V, lo que arroja una relación $V_2/V_1 \approx 1,84$, ligeramente inferior a la relación nominal, lo cual se debe a que la distorsión no lineal aumenta las pérdidas en el núcleo. Esta elevación relativa de la tensión con carga capacitiva (regulación negativa) es consistente con el efecto de la corriente en adelanto sobre la caída de tensión interna del transformador. En este ensayo sí se respeta el balance de potencias ($P_1 = 550$ W $> P_2 = 450$ W), con pérdidas internas razonables. La forma de onda de corriente vs voltaje evidencia la distorsión introducida por el puente rectificador en la corriente demandada por la carga no lineal.

    \item \textbf{Carga inductiva:} En este ensayo la potencia activa del secundario ($P_2 = 624$ W) supera la potencia aparente correspondiente ($S_2 = V_2 \cdot I_2 = 180 \cdot 3 = 540$ VA), lo que implica un factor de potencia superior a la unidad, físicamente imposible. De igual forma, el factor de potencia del primario resulta $\approx 0,2$ ($P_1 = 132$ W sobre $S_1 = 644$ VA), valor poco consistente. Estas discrepancias indican que los vatímetros convencionales, que miden únicamente a la frecuencia fundamental, no son confiables cuando la corriente está fuertemente distorsionada por armónicos generados en el rectificador, por lo que los resultados de potencia en los ensayos no lineales deben interpretarse con cautela. La forma de onda registrada en la carga inductiva muestra la distorsión y el contenido armónico de la corriente esperados en este tipo de carga.
\end{itemize}

En términos generales, los ensayos con carga lineal arrojan resultados más coherentes para el análisis de regulación de tensión y relación de transformación, mientras que las mediciones de potencia en los ensayos con rectificador se ven afectadas por la distorsión armónica y un sistema de medición inadecuado para esas condiciones; por ello dichos valores deben asumirse como aproximados, sujetos a la incertidumbre reportada.